% ============================================================================
% MO DAU (Loi mo dau) — gioi thieu so, dat van de, boi canh
% ============================================================================

\chapter*{MỞ ĐẦU}
\phantomsection
\addcontentsline{toc}{chapter}{MỞ ĐẦU}
\markboth{MỞ ĐẦU}{MỞ ĐẦU}

Trong khoảng một thập kỷ trở lại đây, kiến trúc phần mềm doanh nghiệp đã dịch chuyển mạnh
mẽ từ ứng dụng \emph{một khối} (\textit{monolithic}) sang \emph{vi dịch vụ}
(\textit{microservices}) đóng gói trong container và điều phối bằng Kubernetes. Sự dịch
chuyển này mang lại tốc độ phát triển và khả năng mở rộng vượt trội, nhưng đồng thời làm
thay đổi căn bản bề mặt tấn công: thay vì các lời gọi nội bộ gọn trong một tiến trình, hệ
thống nay có hàng trăm điểm cuối nội bộ liên lạc với nhau qua mạng. Phần lưu lượng nội cụm
(hướng Đông--Tây) này gần như vô hình trước các biện pháp an ninh chu vi truyền thống, trở
thành mảnh đất cho hai hành vi tấn công điển hình: \emph{quét cổng} để trinh sát và
\emph{dịch chuyển ngang} để mở rộng quyền kiểm soát giữa các pod.

\textbf{Đặt vấn đề.} Làm thế nào để \emph{quan sát} được trọn vẹn lưu lượng nội cụm và
\emph{phát hiện sớm} các hành vi quét cổng, dịch chuyển ngang ngay trong một cụm Kubernetes,
với chi phí đủ thấp để chạy thường trực, rồi \emph{khép kín} bằng một cơ chế phản hồi có
kiểm soát? Đây là bài toán mà đồ án hướng tới giải quyết.

\textbf{Nội dung đồ án.} Đồ án xây dựng và đánh giá một hệ thống giám sát --- phản hồi an
ninh đầu cuối cho cụm Kubernetes, đặt tên K8sSecDaP. Hệ thống thu thập sự kiện kết nối ngay
tại \emph{tầng nhân} bằng công nghệ eBPF, phát hiện bất thường bằng các \emph{cấu trúc dữ
liệu luồng} (Count-Min Sketch) và \emph{giải thuật đồ thị} (Tarjan, BFS), quản lý sự cố qua
một cụm dịch vụ SOC, và khép kín vòng phản hồi bằng cách sinh chính sách mạng
\texttt{NetworkPolicy} có người duyệt để cách ly nguồn tấn công. Toàn bộ được triển khai và
đo đạc trên một cụm Kubernetes thật.

\textbf{Lý do chọn đề tài.} Trong bối cảnh nêu trên, mô hình \textit{Zero-Trust} --- nơi
mọi kết nối đều bị coi là không đáng tin cho đến khi được xác thực --- là một nguyên tắc
được giới học thuật và công nghiệp đồng thuận \cite{nist800207}, nhưng việc \textit{hiện
thực hoá} nguyên tắc này trên một cụm Kubernetes nội bộ vẫn còn nhiều vướng mắc kỹ thuật.
Điều khiến Zero-Trust khó hiện thực hóa nội cụm không nằm ở nguyên lý mà ở \emph{khả năng
quan sát}: muốn ``không tin tưởng mặc nhiên mọi kết nối'' thì trước hết phải \emph{thấy}
được mọi kết nối. Đây cũng là lý do đồ án chọn quan sát ở tầng nhân bằng eBPF làm điểm tựa
kỹ thuật: nó cung cấp tầm nhìn Đông--Tây mà các lớp giám sát hướng host không có, từ đó
mới có thể đặt nền cho phát hiện và phản hồi có kiểm soát. Đồ án cũng tiếp nối và mở rộng
kết quả của Báo cáo học phần tốt nghiệp (CNPM~1): từ một bài toán trừu tượng về ước lượng
tần suất trên luồng dữ liệu --- giải bằng Count-Min Sketch --- phát triển thành một hệ
thống hoàn chỉnh từ thu thập ở tầng nhân tới quản lý sự cố và phản hồi trên cụm thật.

\textbf{Bố cục đồ án.} Đồ án gồm năm chương:
\begin{itemize}
  \item \textbf{Chương~\ref{ch:overview} --- Tổng quan:} bối cảnh, khảo sát nhu cầu thực tế, mô hình mối đe dọa, phân tích yêu cầu và khảo sát công cụ liên quan.
  \item \textbf{Chương~\ref{ch:foundations} --- Cơ sở lý thuyết và công nghệ:} eBPF/CO-RE, Count-Min Sketch, Tarjan, BFS, LPM-Trie và nền tảng công nghệ sử dụng.
  \item \textbf{Chương~\ref{ch:design} --- Phân tích và thiết kế hệ thống:} mô hình mối đe dọa chi tiết, kiến trúc phần mềm, thiết kế cơ sở dữ liệu và thiết kế giao diện người dùng.
  \item \textbf{Chương~\ref{ch:impl} --- Cài đặt, triển khai và thử nghiệm hệ thống:} hạ tầng nền tảng, các thành phần lõi (collector/pipeline/SOC), ứng dụng và giao diện vận hành, cùng việc chạy thử nghiệm end-to-end trên cluster thật.
  \item \textbf{Chương~\ref{ch:eval} --- Đánh giá và kết luận:} phương pháp và kết quả đánh giá định lượng, so sánh, giới hạn và hướng phát triển.
\end{itemize}
